\chapter{Landing Experiment Record}
\label{app:landing}

The main text follows L11 as the final step in the landing campaign. This appendix gathers its
identity, configuration and outcome in one place so that the narrative can be checked against the
immutable artifacts. Repository identity remains a submission placeholder until the working tree
is frozen for publication.

\newcommand{\hashvalue}[4]{\begingroup\ttfamily\footnotesize#1\hspace{0pt}#2\hspace{0pt}#3\hspace{0pt}#4\endgroup}
\small
\begin{longtable}{p{0.38\textwidth}p{0.52\textwidth}}
  \toprule
  Field & Value \\
  \midrule
  \endhead
  Run ID and revision & \texttt{L11}, revision 2 (continuation of the same run from 90M to 100M) \\
  Repository commit & TODO after final source freeze \\
  Manifest session identity & \hashvalue{438cf2baccfafe0f}{0a924bf61021a88d}{eb3515c33e74d9d3}{871f539338888deb} \\
  Raw revision-2 \texttt{Session.json} SHA-256 &
  \hashvalue{c166eb67f9b0e810}{b8d40b43b8f0f17f}{33290e4d41119b6a}{bc7bccbf36f45a40} \\
  Model SHA-256 &
  \hashvalue{e03b8facb32a0c40}{ee9449625767d6b4}{315965450028ec3a}{c8f511c8b344830b} \\
  Evaluation-summary SHA-256 &
  \hashvalue{fd48f2dc02351a0c}{752853939ee04b4c}{a267ff0d36df9f59}{796922b3828e329d} \\
  Reward preset/version & Balanced, objective schema 16; success-only mission efficiency \\
  Trainer/environment seeds & 1 / 1 \\
  Training steps and revisions & 100,000,086; two immutable revisions \\
  Final curriculum progress & 100\% linear and effective; 92.90\% recent success \\
  Policy interface & 52 observations; 12 continuous actions; three independent engines \\
  Evaluated model & \texttt{Rocket.onnx}, final checkpoint 100,000,086 \\
  Evaluation seed and episode count & 20257; 250 (50 per fixed difficulty band) \\
  Stable four-foot touchdown rate (95\% CI) & 212/250 = 84.8\% (79.8--88.7\%) \\
  Full-difficulty touchdown rate (95\% CI) & 18/50 = 36\% (24.1--49.9\%) \\
  First-contact speed, vertical/horizontal & \SI{1.38}{m/s}; \SI{1.33}{m/s} / \SI{0.279}{m/s}
  over 248 touchdown episodes \\
  First-contact tilt and angular rate & \SI{0.637}{\degree}; \SI{1.72}{\degree\per\second} \\
  Maximum impulse and rebound & mean maximum \SI{39146}{\newton\second}; mean maximum
  \SI{6.55}{\micro\metre} \\
  Stable-hold time & mean \SI{0.566}{s} \\
  Structural-strike rate & 0/250 = 0\% \\
  Foot-outside-pad rate & 1/250 = 0.4\% \\
  Engine use & mean 3.00 first ignitions and 1.56 restarts per episode \\
  Fuel used & mean \SI{4529}{kg} over all evaluation episodes \\
  Dominant failure modes & 33 hard touchdowns, 3 missed-pad and 2 too-far terminations; 30 hard
  touchdowns occurred at full difficulty \\
  \bottomrule
\end{longtable}
\normalsize

The evaluation-summary file reports \texttt{aborted: false}. No wind, fins, RCS or injected
faults were active. L11 used minimum throttle 0.499, startup delay \SI{0.45}{s}, shutdown
transient \SI{0.20}{s}, minimum run \SI{1.5}{s}, cooldown \SI{2.0}{s} and throttle spool rate
5.8 per second. Its success-only efficiency magnitude, restart equivalent and touchdown-quality
fraction were increased relative to L10 to favour longer useful burns without making early
shutdown attractive on failed episodes.

The full-difficulty result is based on a model trained at 100\% curriculum, but it remains much
lower than exact-$d=1$ training success. This discrepancy is retained as a finding rather than
removed through checkpoint selection. The next controlled diagnostic is the deterministic versus
stochastic inference comparison described in Chapter~\ref{ch:limitations}.

Before submission, add representative seed-labelled success and failure videos and the final
public repository commit. Video remains qualitative evidence; the numeric result above comes
from the complete evaluation summary.
